\documentclass[11pt,a4paper]{article}

% ============================================================
% PAQUETES
% ============================================================
\usepackage[utf8]{inputenc}
\usepackage[T1]{fontenc}
\usepackage[spanish,es-tabla]{babel}
\usepackage{lmodern}
\usepackage{microtype}
\usepackage[a4paper,margin=2.2cm,headheight=15pt]{geometry}
\usepackage{amsmath,amssymb}
\usepackage{booktabs,longtable,tabularx,array}
\usepackage[table]{xcolor}
\usepackage{enumitem}
\usepackage{float}
\usepackage{titlesec}
\usepackage{fancyhdr}
\usepackage{lastpage}
\usepackage{hyperref}
\usepackage{tcolorbox}
\usepackage{underscore}
\usepackage{ragged2e}

\tcbuselibrary{breakable,skins}

% ============================================================
% COLORES Y ESTILOS
% ============================================================
\definecolor{azul}{HTML}{1F3864}
\definecolor{azulclaro}{HTML}{D9E2F3}
\definecolor{verde}{HTML}{2E7D32}
\definecolor{naranja}{HTML}{C55A11}
\definecolor{gris}{HTML}{F2F2F2}
\definecolor{grisoscuro}{HTML}{595959}

\newcommand{\code}[1]{\texttt{#1}}
\newcolumntype{Y}{>{\RaggedRight\arraybackslash}X}
\newcolumntype{L}[1]{>{\RaggedRight\arraybackslash}p{#1}}

\tcbset{
  idea/.style={colback=azul!5,colframe=azul,fonttitle=\bfseries,breakable,
    left=4pt,right=4pt,top=4pt,bottom=4pt},
  ojo/.style={colback=naranja!5,colframe=naranja,fonttitle=\bfseries,breakable,
    left=4pt,right=4pt,top=4pt,bottom=4pt},
  frase/.style={colback=verde!5,colframe=verde,fonttitle=\bfseries,breakable,
    left=4pt,right=4pt,top=4pt,bottom=4pt},
  dato/.style={colback=gris,colframe=grisoscuro,fonttitle=\bfseries,breakable,
    left=4pt,right=4pt,top=4pt,bottom=4pt}
}

\titleformat{\section}
  {\Large\bfseries\color{azul}}{\thesection}{0.6em}{}
  [\vspace{-0.25em}{\color{azul}\rule{\textwidth}{0.8pt}}]
\titleformat{\subsection}
  {\large\bfseries\color{azul}}{\thesubsection}{0.5em}{}
\titleformat{\subsubsection}
  {\normalsize\bfseries\color{azul}}{\thesubsubsection}{0.5em}{}
\titlespacing*{\section}{0pt}{1.2ex plus .2ex}{0.8ex}
\titlespacing*{\subsection}{0pt}{1.0ex plus .2ex}{0.5ex}

\pagestyle{fancy}
\fancyhf{}
\fancyhead[L]{\small\color{grisoscuro}Guía de estudio — Fase 2}
\fancyhead[R]{\small\color{grisoscuro}Sistema de Brigadas de Salud}
\fancyfoot[C]{\small\thepage\ de \pageref{LastPage}}
\renewcommand{\headrulewidth}{0.4pt}

\hypersetup{
  colorlinks=true,
  linkcolor=azul,
  urlcolor=azul,
  citecolor=azul,
  pdftitle={Guía de estudio y defensa — Fase 2},
  pdfauthor={Jorman David Noriega Dávila; Julio César Ríos Maldonado; Gabriel Omar Venegas González}
}

\setlength{\parskip}{0.45em}
\setlength{\parindent}{0pt}
\renewcommand{\arraystretch}{1.18}

% ============================================================
% DOCUMENTO
% ============================================================
\begin{document}

% ---------- PORTADA ----------
\hypersetup{pageanchor=false}
\begin{titlepage}
\centering
\vspace*{1.2cm}

{\color{azul}\rule{\textwidth}{3pt}}\\[0.45cm]
{\Huge\bfseries\color{azul} Guía de estudio\\[0.15cm] y defensa — Fase 2}\\[0.55cm]
{\Large Sistema Inteligente de Priorización y\\[0.1cm]
Optimización Espacial de Brigadas de Salud Extramurales}\\[0.35cm]
{\color{azul}\rule{\textwidth}{1pt}}\\[1.0cm]

\begin{tabular}{r l}
\textbf{Equipo:} & Jorman David Noriega Dávila \\
 & Julio César Ríos Maldonado \\
 & Gabriel Omar Venegas González \\[0.35cm]
\textbf{Docente:} & Tonny Jiménez \\
\textbf{Universidad:} & Universidad Popular del Cesar \\
\textbf{Facultad:} & Ingenierías y Tecnologías \\
\textbf{Programa:} & Ingeniería de Sistemas e Informática \\
\textbf{Ciudad:} & Valledupar, Cesar \\
\textbf{Año:} & 2026 \\
\end{tabular}

\vfill

\begin{tcolorbox}[colback=azul!5,colframe=azul,width=0.88\textwidth]
\centering
\small
Documento de apoyo para la exposición.\\
No reemplaza el documento técnico de 12 páginas de la Fase 2.
\end{tcolorbox}

\vspace{0.6cm}
\end{titlepage}
\hypersetup{pageanchor=true}

% ---------- RESUMEN ----------
\begin{abstract}
\noindent
Esta guía reúne, en un solo documento, la explicación del dataset que dejó la Fase~1, el procesamiento y los resultados de la Fase~2, y un plan de estudio para explicar el proyecto paso a paso. Está organizada en tres bloques: \textbf{(1)} qué contiene la tabla \code{municipio_anio.parquet}, cuáles son sus 27 entradas y cuál es su salida; \textbf{(2)} qué hacen los modelos y el MILP con esa información; y \textbf{(3)} cómo estudiar y exponer el proyecto, con guion de seis minutos, banco de preguntas y errores que conviene evitar.
\end{abstract}

\begingroup
\small
\tableofcontents
\endgroup
\newpage

% ============================================================
\section{Cómo usar esta guía}
% ============================================================

Esta guía está pensada para tres momentos:

\begin{enumerate}[leftmargin=*]
  \item \textbf{Entender:} leer las secciones 1 y 2 para dominar el dataset y el modelado.
  \item \textbf{Estudiar:} seguir la ruta de tres días de la sección 3.
  \item \textbf{Exponer:} ensayar el guion de la sección 4 y repasar el banco de preguntas de la sección 5.
\end{enumerate}

Si tienes poco tiempo, prioriza: \textbf{1.3} (target), \textbf{2.3} (métricas), \textbf{2.5} (MILP), \textbf{4} (guion) y \textbf{5} (preguntas).

\subsection{Resumen ejecutivo en 12 números}

\begin{table}[H]
\centering
\small
\begin{tabularx}{\textwidth}{@{}c L{7.2cm} Y@{}}
\toprule
\textbf{\#} & \textbf{Dato} & \textbf{Valor} \\
\midrule
1 & Municipios en el dataset & 1.122 \\
2 & Filas municipio--año & 5.610 \\
3 & Columnas totales & 74 \\
4 & Variables de entrada del modelo & 27 \\
5 & Filas usables para modelar & 5.394 \\
6 & Positivos (\code{prioridad_alta}) & 1.798 (33,3\%) \\
7 & Train / Validación / Test (filas) & 3.770 / 807 / 817 \\
8 & Train / Validación / Test (municipios) & 785 / 168 / 169 \\
9 & Mejor modelo & HistGradientBoosting \\
10 & F1 / recall / ROC-AUC del mejor modelo & 0,6234 / 0,5379 / 0,7848 \\
11 & Brigadas asignadas por el MILP & 4 \\
12 & Municipios del Cesar evaluados & 25 \\
\bottomrule
\end{tabularx}
\end{table}

\subsection{Ruta de estudio de tres días}

\begin{table}[H]
\centering
\small
\begin{tabularx}{\textwidth}{@{}l c c Y@{}}
\toprule
\textbf{Día} & \textbf{Bloques} & \textbf{Tiempo} & \textbf{Producto esperado} \\
\midrule
Día 1 & 1 y 2 & 2,5 h & Explicar la tabla y la fórmula del target en una hoja. \\
Día 2 & 3 y 4 & 3,0 h & Explicar los modelos y recalcular la matriz de confusión. \\
Día 3 & 5 y 6 & 2,5 h & Explicar el MILP y ensayar el guion de seis minutos. \\
\bottomrule
\end{tabularx}
\end{table}

% ============================================================
\section{Qué dejó la Fase 1: el dataset}
% ============================================================

\subsection{Resumen en una frase}

La Fase~1 no entregó un modelo: entregó una \textbf{tabla analítica municipio--año} que cruza 1.122 municipios de Colombia durante 2017--2021 con datos de RIPS, REPS, NBI, resguardos, población y DIVIPOLA. Esa tabla es la materia prima de la Fase~2.

\subsection{Ficha técnica}

\begin{table}[H]
\centering
\small
\begin{tabularx}{\textwidth}{@{}L{4.2cm} Y@{}}
\toprule
\textbf{Elemento} & \textbf{Valor} \\
\midrule
Archivo principal & \code{fase 1 api/data/modelo/municipio_anio.parquet} \\
Formato alterno & \code{municipio_anio.csv} \\
Unidad de análisis & Municipio $\times$ año \\
Llave & \code{cod_mpio} + \code{a_o} \\
Filas & 5.610 (1.122 municipios $\times$ 5 años) \\
Columnas & 74 \\
Años & 2017--2021 \\
Filas usables para modelar & 5.394 \\
Entradas del modelo & 27 \\
Salida del modelo & \code{prioridad_alta} (binaria) \\
Train / Validación / Test (filas) & 3.770 / 807 / 817 \\
Train / Validación / Test (municipios) & 785 / 168 / 169 \\
Positivos & 1.798 de 5.394 = 33,3\% \\
\bottomrule
\end{tabularx}
\end{table}

\begin{tcolorbox}[ojo,title={La advertencia más importante}]
No son 5.610 municipios independientes. Son \textbf{1.122 municipios con 5 observaciones cada uno}. Además, muchas variables estructurales ---NBI, REPS, resguardos y los porcentajes CIE-10--- son prácticamente constantes en los cinco años. La variación anual viene sobre todo de la población y del volumen de atenciones. Por eso el modelo es \textbf{transversal}, no una serie de tiempo, y el split se hace por municipio, no por fila.
\end{tcolorbox}

\subsection{La salida: \code{prioridad_alta}}

No es ``el municipio tiene o no déficit'' en términos absolutos. Es una \textbf{prioridad relativa} construida así:

\begin{align}
\text{tasa}_{i,t} &= \frac{\text{consultas}_{i,t}}{\text{población}_{i,t}} \\
\text{bench}_{i,t} &= \operatorname{mediana}\left\{\text{tasa}_{j,t} : j \in \text{banda}_i\right\} \\
\text{brecha}_{i,t} &= \max\left(0,\; 1 - \frac{\text{tasa}_{i,t}}{\text{bench}_{i,t}}\right) \\
\text{prioridad alta}_{i,t} &=
\begin{cases}
1 & \text{si } \text{brecha}_{i,t} \ge q_{66} \approx 0{,}337 \\
0 & \text{en otro caso}
\end{cases}
\end{align}

\paragraph{Cómo se interpreta.}
\begin{itemize}[leftmargin=*]
  \item Se compara cada municipio con municipios \textbf{de tamaño similar}, no contra el promedio nacional.
  \item Si un municipio tiene una tasa de consultas muy por debajo de sus pares, su \code{brecha} sube.
  \item El corte está en el \textbf{tercio superior} de esa brecha: aproximadamente \textbf{0,337}.
  \item Eso significa: ``este municipio tiene cerca de \textbf{34\% menos consultas por habitante} que la mediana de municipios parecidos''.
  \item El \textbf{33,3\% de positivos es por construcción}, no porque el 33\% de Colombia tenga un déficit crítico. Es una etiqueta relativa de priorización.
\end{itemize}

\paragraph{Ejemplos reales con datos de 2021.}

\begin{table}[H]
\centering
\small
\begin{tabular}{@{}l r r r c@{}}
\toprule
\textbf{Municipio} & \textbf{Consultas/1.000 hab} & \textbf{Benchmark} & \textbf{Brecha} & \textbf{Etiqueta} \\
\midrule
Valledupar & 2.388 & 1.903 & 0,000 & 0 \\
Pueblo Bello & 2.679 & 1.492 & 0,000 & 0 \\
Becerril & 1.195 & 1.492 & 0,199 & 0 \\
Astrea & 728 & 1.492 & 0,512 & 1 \\
Chimichagua & 873 & 1.903 & 0,541 & 1 \\
\bottomrule
\end{tabular}
\end{table}

Esto explica algo que suele confundir: \textbf{Pueblo Bello no es ``prioridad alta'' para el modelo} porque su tasa de consultas está por encima de la mediana de sus pares. Luego veremos por qué aun así el MILP lo selecciona.

\paragraph{¿Qué es un hub y por qué se excluye?}
Un \textbf{hub} es un prestador de referencia regional: su tasa de consultas es al menos 3 veces la mediana de su banda de población. Atrae pacientes de otros municipios, así que si se incluyera en el ranking aparecería como ``muy bien cubierto'' y distorsionaría la comparación.

\begin{itemize}[leftmargin=*]
  \item Hay \textbf{211 filas municipio--año} marcadas como hub.
  \item Se conservan en la tabla, pero \textbf{no entran al entrenamiento}.
  \item Del total, quedan \textbf{5.394 filas usables} (211 hubs y 5 filas sin RIPS que ya no tenían target).
\end{itemize}

\subsection{Las 27 entradas del modelo}

El dataset tiene 74 columnas, pero el modelo usa solo \textbf{27 features}. Están agrupadas en cinco bloques:

\begin{table}[H]
\centering
\small
\begin{tabularx}{\textwidth}{@{}l c Y@{}}
\toprule
\textbf{Bloque} & \textbf{\#} & \textbf{Variables y qué representan} \\
\midrule
Geografía y demografía & 5 &
\code{latitud}, \code{longitud}, \code{poblacion}, \code{pct_rural_disperso}, \code{banda_pob}.
Ubicación, tamaño y ruralidad. \\
Pobreza multidimensional (NBI) & 11 &
\code{nbi_total}, \code{nbi_miseria}, \code{nbi_vivienda}, \code{nbi_servicios},
\code{nbi_hacinamiento}, \code{nbi_inasistencia}, \code{nbi_dependencia},
\code{nbi_rural}, \code{nbi_rural_miseria}, \code{nbi_rural_servicios},
\code{nbi_rural_hacinamiento}. Vulnerabilidad estructural y pobreza rural. \\
Capacidad instalada (REPS) & 3 &
\code{n_sedes}, \code{n_sedes_ese}, \code{sedes_10k_hab}.
Oferta de servicios de salud. \\
Contexto étnico & 3 &
\code{n_resguardos}, \code{pob_resguardos}, \code{n_comunidades}.
Presencia indígena y equidad. \\
Perfil de morbilidad (CIE-10) & 5 &
\code{pct_cap_O}, \code{pct_cap_A}, \code{pct_cap_J}, \code{pct_cap_K}, \code{pct_cap_Z}.
Mezcla de diagnósticos. \\
\bottomrule
\end{tabularx}
\end{table}

\paragraph{Qué significa cada grupo.}
\begin{itemize}[leftmargin=*]
  \item \textbf{Geografía y demografía:} la hipótesis es que municipios rurales, dispersos y pequeños tienen más barreras de acceso.
  \item \textbf{NBI:} mide pobreza multidimensional. \code{nbi_servicios}, por ejemplo, es el porcentaje de viviendas sin servicios adecuados; \code{nbi_rural_servicios} es lo mismo en el área rural dispersa.
  \item \textbf{REPS:} mide capacidad instalada. \code{sedes_10k_hab} es el número de sedes por cada 10.000 habitantes.
  \item \textbf{Étnico:} \code{pob_resguardos} es la población en resguardos indígenas; \code{n_resguardos} cuenta cuántos hay.
  \item \textbf{CIE-10:} son \textbf{porcentajes} de atenciones por capítulo diagnóstico, no conteos.
\end{itemize}

\paragraph{Los cinco capítulos CIE-10 usados.}

\begin{table}[H]
\centering
\small
\begin{tabularx}{\textwidth}{@{}c L{5.0cm} Y@{}}
\toprule
\textbf{Código} & \textbf{Capítulo} & \textbf{Interpretación} \\
\midrule
\code{pct_cap_O} & Embarazo, parto y puerperio & Demanda materno-infantil. \\
\code{pct_cap_A} & Enfermedades infecciosas y parasitarias & Carga infecciosa. \\
\code{pct_cap_J} & Enfermedades respiratorias & Carga respiratoria. \\
\code{pct_cap_K} & Enfermedades digestivas & Carga digestiva. \\
\code{pct_cap_Z} & Contacto con servicios de salud & Utilización de servicios; no es una enfermedad. \\
\bottomrule
\end{tabularx}
\end{table}

El dataset tiene 26 capítulos CIE-10 (A--Z). El modelo usa cinco porque concentran la demanda más relevante para brigadas extramurales: materno-infantil, infecciosas, respiratorias, digestivas y contacto con servicios. Los otros 21 quedaron fuera para reducir ruido, redundancia y complejidad. No porque no existan.

\subsection{Lo que está en el dataset pero \textbf{no} entra al modelo}

Esto es clave para defender que no hay fuga de datos directa:

\begin{table}[H]
\centering
\footnotesize
\begin{tabularx}{\textwidth}{@{}L{4.2cm} Y L{4.3cm}@{}}
\toprule
\textbf{Grupo} & \textbf{Columnas} & \textbf{Por qué se excluyen} \\
\midrule
Volumen y tasas RIPS &
\code{consultas}, \code{hospitalizaciones}, \code{procedimientos}, \code{urgencias},
\code{atenciones_total}, \code{consultas_1000hab}, \code{urgencias_1000hab},
\code{hosp_1000hab}, \code{proced_1000hab}, \code{ratio_consultas_urgencias} &
Son la materia prima del target; incluirlas sería fuga. \\
Derivadas del target &
\code{bench_consultas}, \code{brecha_consultas}, \code{bench_global}, \code{brecha_global},
\code{brecha_nacional}, \code{prioridad_alta}, \code{es_hub} &
Contienen la respuesta. \\
Identificadores y control &
\code{cod_mpio}, \code{cod_dpto}, \code{nom_mpio}, \code{a_o}, \code{split}, \code{sin_rips} &
No son predictores. \\
CIE-10 restantes & 21 capítulos & Redundancia y ruido. \\
NBI restantes &
\code{nbi_rural_vivienda}, \code{nbi_rural_inasistencia}, \code{nbi_rural_dependencia} &
Redundancia con las otras NBI. \\
\bottomrule
\end{tabularx}
\end{table}

La regla es: \textbf{el modelo usa contexto estructural y perfil de morbilidad; no usa el volumen ni la brecha de consultas}. Eso es lo que hace que la predicción tenga sentido.

\subsection{Las tres ideas que debes dominar de la Fase 1}

\begin{tcolorbox}[idea,title={Ideas clave}]
\begin{enumerate}[leftmargin=*,itemsep=2pt]
  \item \textbf{La tabla es municipio--año, no paciente.} 5.610 filas, 74 columnas, 1.122 municipios, 2017--2021.
  \item \textbf{El target es relativo:} tercio superior de la brecha contra municipios pares. No es prevalencia real.
  \item \textbf{El split es por municipio:} 785 municipios en train, 168 en validación y 169 en test. Ningún municipio aparece en dos conjuntos.
\end{enumerate}
\end{tcolorbox}

% ============================================================
\section{Qué hace la Fase 2 con esos datos}
% ============================================================

\subsection{Pipeline en siete pasos}

\begin{table}[H]
\centering
\small
\begin{tabularx}{\textwidth}{@{}c L{4.2cm} Y@{}}
\toprule
\textbf{Paso} & \textbf{Acción} & \textbf{Resultado} \\
\midrule
1 & Cargar y limpiar & Se excluyen hubs, se eliminan filas sin target y se imputan medianas de CIE-10. Quedan 5.394 filas. \\
2 & Separar & Train, validación y test vienen definidos desde Fase~1, agrupados por municipio. \\
3 & Escalar & Solo para Regresión Logística, con \code{StandardScaler}. Los árboles no necesitan escalado. \\
4 & Entrenar & Cuatro modelos: Dummy, Regresión Logística, Random Forest e HistGradientBoosting. \\
5 & Evaluar & Métricas en test ciego: accuracy, precision, recall, F1, ROC-AUC, PR-AUC y matriz de confusión. \\
6 & Interpretar & Importancia Gini del Random Forest. \\
7 & Optimizar & La probabilidad del mejor modelo alimenta el MILP de asignación de brigadas. \\
\bottomrule
\end{tabularx}
\end{table}

\subsection{Modelos y ajuste}

\begin{table}[H]
\centering
\small
\begin{tabularx}{\textwidth}{@{}L{3.6cm} L{2.4cm} c Y@{}}
\toprule
\textbf{Modelo} & \textbf{Tipo} & \textbf{Escalado} & \textbf{Ajuste y rol} \\
\midrule
Dummy & Baseline trivial & No &
\code{most_frequent}: predice siempre la clase mayoritaria. Sirve para demostrar por qué la accuracy sola engaña. \\
Regresión Logística & Lineal & Sí &
\code{class_weight='balanced'}, \code{max_iter=1000}. Modelo interpretable y base lineal. \\
Random Forest & Bagging de árboles & No &
Grid de 12 combinaciones $\times$ 3 folds: \code{n_estimators} [100, 200], \code{max_depth} [6, 10, None], \code{min_samples_split} [2, 5], \code{class_weight='balanced'}. \\
HistGradientBoosting & Boosting de árboles & No &
Grid de 24 combinaciones $\times$ 3 folds: \code{learning_rate} [0.03, 0.08, 0.15], \code{max_iter} [100, 150], \code{max_leaf_nodes} [15, 31], \code{min_samples_leaf} [20, 40]. \\
\bottomrule
\end{tabularx}
\end{table}

\begin{itemize}[leftmargin=*]
  \item La búsqueda de hiperparámetros usa \code{GridSearchCV} con \code{cv=3} y \code{scoring='f1'}.
  \item El escalado se aplica solo a la Regresión Logística; los árboles son invariantes a la escala.
  \item \code{class_weight='balanced'} compensa que la clase positiva sea aproximadamente un tercio del total.
  \item El conjunto de test es ciego: no se usa para elegir hiperparámetros ni umbral.
\end{itemize}

\subsection{Resultados reales en test}

Test: \textbf{817 filas, 277 positivos (33,9\%)}.

\begin{table}[H]
\centering
\small
\begin{tabular}{@{}l r r r r r r@{}}
\toprule
\textbf{Modelo} & \textbf{Accuracy} & \textbf{Precision} & \textbf{Recall} & \textbf{F1} & \textbf{ROC-AUC} & \textbf{PR-AUC} \\
\midrule
\textbf{HistGradientBoosting} & \textbf{0,7797} & \textbf{0,7413} & \textbf{0,5379} & \textbf{0,6234} & \textbf{0,7848} & \textbf{0,7414} \\
Random Forest & 0,7283 & 0,6211 & 0,5090 & 0,5595 & 0,7653 & 0,6803 \\
Regresión Logística & 0,6573 & 0,4952 & 0,5596 & 0,5254 & 0,6920 & 0,5926 \\
Dummy & 0,6610 & 0,0000 & 0,0000 & 0,0000 & 0,5000 & 0,3390 \\
\bottomrule
\end{tabular}
\end{table}

\paragraph{Matriz de confusión de HistGradientBoosting.}

\begin{table}[H]
\centering
\small
\begin{tabular}{@{}l c c@{}}
\toprule
 & \textbf{Predice no prioritario} & \textbf{Predice prioritario} \\
\midrule
\textbf{Real no prioritario} & 488 TN & 52 FP \\
\textbf{Real prioritario} & 128 FN & 149 TP \\
\bottomrule
\end{tabular}
\end{table}

\paragraph{Cómo leer esto.}
\begin{itemize}[leftmargin=*]
  \item \textbf{Dummy tiene 66,1\% de accuracy y F1 = 0.} ¿Por qué? Porque el 66,1\% del test son negativos; si predice siempre ``no prioritario'', acierta la mayoría, pero \textbf{no detecta ni un solo municipio prioritario}. Por eso la accuracy sola engaña.
  \item \textbf{HGB es el mejor en F1 y PR-AUC.} Detecta \textbf{149 de 277} municipios prioritarios: recall \textbf{53,8\%}. Cuando alerta, acierta el \textbf{74,1\%} de las veces: precision \textbf{0,7413}.
  \item \textbf{ROC-AUC 0,7848:} si tomas un municipio prioritario y uno no prioritario al azar, el modelo ordena al prioritario por encima en el 78,5\% de los casos.
  \item \textbf{PR-AUC 0,7414:} mucho mejor que el baseline de 0,339 (la prevalencia). Es la métrica más informativa cuando la clase positiva es minoritaria.
  \item \textbf{El umbral es 0,5.} No es sagrado: si el objetivo es no dejar municipios por fuera, se puede bajar el umbral y ganar recall a costa de precisión. Eso es una decisión operativa.
\end{itemize}

\begin{tcolorbox}[idea,title={Cómo explicar la matriz en una frase}]
``El modelo detecta 149 de los 277 municipios prioritarios; se le escapan 128 y hace 52 falsas alarmas. Cuando dice `prioritario', acierta tres de cada cuatro veces.''
\end{tcolorbox}

\subsection{Interpretabilidad: top 5 del Random Forest}

\begin{table}[H]
\centering
\small
\begin{tabular}{@{}l r@{}}
\toprule
\textbf{Variable} & \textbf{Importancia Gini} \\
\midrule
\code{pct_cap_O} (materno) & 0,11218 \\
\code{pob_resguardos} & 0,09068 \\
\code{nbi_servicios} & 0,08490 \\
\code{nbi_rural_servicios} & 0,05946 \\
\code{pct_cap_A} (infecciosas) & 0,04852 \\
\bottomrule
\end{tabular}
\end{table}

Por bloques, aproximadamente:

\begin{table}[H]
\centering
\small
\begin{tabular}{@{}l r@{}}
\toprule
\textbf{Bloque} & \textbf{Importancia aproximada} \\
\midrule
NBI & 42\% \\
CIE-10 & 26\% \\
Geografía y demografía & 15\% \\
Étnico & 12\% \\
REPS & 5\% \\
\bottomrule
\end{tabular}
\end{table}

\begin{tcolorbox}[ojo,title={Advertencia sobre Gini}]
Gini mide reducción de impureza, \textbf{no causalidad}. Si dos variables están correlacionadas, se reparten la importancia. \code{pct_cap_O} puede ser un proxy de demanda materna, no una causa. No digas ``esta variable causa la prioridad''; di ``esta variable es la más asociada en el modelo''.
\end{tcolorbox}

\subsection{El MILP: de probabilidad a plan de brigadas}

El modelo anterior puntúa municipios, pero no decide un plan. El MILP sí:

\begin{equation}
\max \sum_{i \in \mathcal{M}}
\left[\, p_i \cdot \log(1 + v_i) - \lambda d_i \,\right] x_i
\end{equation}

sujeto a:

\begin{align}
\sum_{i \in \mathcal{M}} x_i &\le K \\
\sum_{i \in \mathcal{R}} x_i &\ge 2 \\
x_i &= 0 \quad \text{si } d_i > D_{\max} \\
x_i &\in \{0,1\}
\end{align}

Donde:

\begin{table}[H]
\centering
\small
\begin{tabularx}{\textwidth}{@{}L{3.2cm} Y@{}}
\toprule
\textbf{Símbolo} & \textbf{Significado} \\
\midrule
$p_i$ & Probabilidad de prioridad alta del municipio $i$, estimada por el HGB. \\
$v_i$ & Población vulnerable: población rural dispersa + población en resguardos. \\
$d_i$ & Distancia Haversine desde Valledupar (hospital base). \\
$\lambda$ & Penalización por distancia: 0,005. \\
$K$ & Brigadas disponibles: 4. \\
$D_{\max}$ & Distancia máxima por brigada: 350 km. \\
$\mathcal{R}$ & Municipios con al menos un resguardo indígena. \\
\bottomrule
\end{tabularx}
\end{table}

\paragraph{Resultado en el Cesar, 2021.}

\begin{table}[H]
\centering
\small
\begin{tabular}{@{}l r r r c@{}}
\toprule
\textbf{Municipio} & \textbf{Prob. IA} & \textbf{Distancia} & \textbf{Resguardos} & \textbf{Brigada} \\
\midrule
Pueblo Bello & 36,9\% & 36,06 km & 2 & \#1 \\
Becerril & 45,8\% & 84,10 km & 2 & \#2 \\
Astrea & 90,9\% & 132,70 km & 0 & \#3 \\
Chimichagua & 75,0\% & 146,76 km & 0 & \#4 \\
\bottomrule
\end{tabular}
\end{table}

\paragraph{Por qué Pueblo Bello entra con 36,9\%.}
Porque el MILP no solo maximiza probabilidad. También exige \textbf{mínimo dos municipios con resguardos}. Pueblo Bello y Becerril cumplen esa cuota. No es una inconsistencia: son dos criterios distintos.

\begin{itemize}[leftmargin=*]
  \item El \textbf{modelo} dice: ``la cobertura de Pueblo Bello no está críticamente por debajo de sus pares''.
  \item El \textbf{MILP} dice: ``aun así, por equidad, debe entrar un municipio con resguardos''.
\end{itemize}

\begin{tcolorbox}[frase,title={Pregunta típica}]
\textbf{``¿Por qué el modelo dice 36,9\% y aun así lo asignan?''}\\
Porque el MILP tiene una restricción de equidad: mínimo dos municipios con resguardos. El modelo mide brecha de cobertura; el MILP agrega equidad, distancia y población vulnerable. Son criterios complementarios.
\end{tcolorbox}

\begin{tcolorbox}[ojo,title={No es ruteo}]
El MILP es un modelo de \textbf{asignación de brigadas a municipios}, no un problema de ruteo de vehículos (VRP). La distancia es Haversine desde la base, no una ruta por carretera. El ruteo real es trabajo futuro.
\end{tcolorbox}

\subsection{La demo}

\begin{verbatim}
python "fase 2/src/probar_ia.py" "Valledupar"
python "fase 2/src/probar_ia.py" "Pueblo Bello"
python "fase 2/src/probar_ia.py" "Astrea"
python "fase 2/src/optimizacion_milp.py"
\end{verbatim}

\code{probar_ia.py} carga el parquet y el modelo HGB ya entrenados, filtra el municipio en 2021 y muestra probabilidad y recomendación. \textbf{No reentrena.}

\paragraph{Qué esperar.}
\begin{itemize}[leftmargin=*]
  \item Valledupar: probabilidad muy baja ($\approx$ 0,9\%), recomendación normal.
  \item Pueblo Bello: 36,9\%, recomendación normal para el modelo, pero entra al plan del MILP por la cuota de resguardos.
  \item Astrea: 90,9\%, prioridad alta.
\end{itemize}

% ============================================================
\section{Plan de estudio}
% ============================================================

\subsection{Ruta de tres días}

\begin{table}[H]
\centering
\small
\begin{tabularx}{\textwidth}{@{}l c c Y@{}}
\toprule
\textbf{Día} & \textbf{Bloques} & \textbf{Tiempo} & \textbf{Producto} \\
\midrule
Día 1 & 1 y 2 & 2,5 h & Explicar la tabla y la fórmula en una hoja. \\
Día 2 & 3 y 4 & 3,0 h & Resolver la matriz de confusión a mano. \\
Día 3 & 5 y 6 & 2,5 h & Guion de seis minutos sin notas. \\
\bottomrule
\end{tabularx}
\end{table}

Si tienes una semana, haz un bloque por día y repite el ensayo dos veces.

\subsection{Bloque 1 --- El dataset y el target (90 min)}

\paragraph{Estudia.}
\begin{itemize}[leftmargin=*]
  \item \code{fase 1 api/fase 1/docs/diccionario_datos.md} (completo; es el mejor resumen).
  \item \code{fase 1 api/fase 1/src/build_tabla_modelo.py}, líneas 120--154 (target) y 38--56 (split).
\end{itemize}

\paragraph{Debes poder explicar.}
\begin{itemize}[leftmargin=*]
  \item Qué es una fila: municipio--año.
  \item Por qué 5.610 = 1.122 $\times$ 5.
  \item La fórmula del target y el corte 0,337.
  \item Por qué 33,3\% de positivos es por construcción.
  \item Qué es un hub y por qué se excluye.
  \item Por qué el split es por municipio.
\end{itemize}

\paragraph{Ejercicio.}
Con los datos de la tabla de ejemplos, calcula a mano la brecha de Astrea y de Valledupar. Verifica que Astrea queda en 1 y Valledupar en 0.

\begin{tcolorbox}[frase,title={Frase para el profesor}]
``La Fase~1 produjo una tabla municipio--año de 5.610 filas y 74 columnas. La salida es una etiqueta relativa: el tercio de municipios con mayor brecha de consultas frente a municipios de tamaño similar.''
\end{tcolorbox}

\subsection{Bloque 2 --- Features y fuga de datos (60--90 min)}

\paragraph{Estudia.}
\begin{itemize}[leftmargin=*]
  \item \code{fase 1 api/fase 2/src/modelado.py}, líneas 36--69 (lista \code{FEATURES}).
  \item \code{fase 1 api/fase 1/docs/diccionario_datos.md}, sección ``Advertencias''.
\end{itemize}

\paragraph{Debes poder explicar.}
\begin{itemize}[leftmargin=*]
  \item 74 columnas del dataset vs 27 entradas del modelo.
  \item Los cinco bloques de variables.
  \item Por qué \code{consultas_1000hab}, \code{brecha_consultas} y \code{prioridad_alta} no son entradas.
  \item Por qué los CIE-10 son porcentajes y no conteos.
  \item Que los CIE-10, NBI, REPS y resguardos son casi constantes en el tiempo; el modelo es transversal.
\end{itemize}

\paragraph{Pregunta difícil.}
``Si el target sale de RIPS, ¿usar CIE-10 no es fuga?''\\
\textbf{Respuesta:} ``El target es volumen/tasa de consultas; las features CIE-10 son la composición porcentual de diagnósticos, no el conteo. Excluimos volumen, tasas, benchmarks y brechas. Aun así, \code{pct_cap_Z} es la variable más discutible porque mide contacto con servicios; en una versión estricta la quitaría y volvería a evaluar.''

\subsection{Bloque 3 --- Modelos y diseño experimental (90 min)}

\paragraph{Estudia.}
\begin{itemize}[leftmargin=*]
  \item \code{fase 1 api/fase 2/src/modelado.py}, líneas 130--244.
  \item \code{fase 1 api/fase 2/reportes/metricas_comparacion.csv}.
\end{itemize}

\paragraph{Debes poder explicar.}
\begin{itemize}[leftmargin=*]
  \item Dummy: baseline que predice la clase mayoritaria.
  \item Regresión Logística: lineal, interpretable, necesita escalado.
  \item Random Forest: muchos árboles, bagging, no necesita escalado.
  \item HistGradientBoosting: boosting, árboles secuenciales, captura no linealidades.
  \item GridSearch: RF 12 combinaciones $\times$ 3 folds; HGB 24 combinaciones $\times$ 3 folds.
  \item Por qué se usa \code{class_weight='balanced'} en LR y RF.
  \item Por qué el test es ciego.
\end{itemize}

\paragraph{Pregunta difícil.}
``¿Por qué no una red neuronal?''\\
\textbf{Respuesta:} ``Con 3.770 filas de entrenamiento y 27 variables tabulares, los ensambles de árboles son más robustos, rápidos e interpretables. Una red neuronal requeriría más datos y tuning, y no garantiza mejor generalización.''

\paragraph{Pregunta difícil.}
``¿Usaron el set de validación?''\\
\textbf{Respuesta honesta:} ``El split 70/15/15 está definido y el set de validación está reservado, pero en la versión actual la selección de hiperparámetros se hizo con validación cruzada de 3 pliegues sobre train y la evaluación final con test. El set de validación queda para calibrar el umbral y para la siguiente iteración.''

\subsection{Bloque 4 --- Métricas (90 min)}

\paragraph{Estudia.}
\begin{itemize}[leftmargin=*]
  \item La tabla de métricas y la matriz de confusión.
  \item \code{fase 1 api/fase 2/reportes/matriz_confusion_modelos.png}.
  \item \code{fase 1 api/fase 2/reportes/curva_roc_comparacion.png}.
\end{itemize}

\paragraph{Debes poder calcular a mano.}
Para HGB:
\begin{align*}
\text{Precision} &= \frac{149}{149 + 52} = 0{,}7413 \\
\text{Recall} &= \frac{149}{149 + 128} = 0{,}5379 \\
\text{F1} &= \frac{2 \times 0{,}7413 \times 0{,}5379}{0{,}7413 + 0{,}5379} = 0{,}6234 \\
\text{Accuracy} &= \frac{488 + 149}{817} = 0{,}7797
\end{align*}

\paragraph{Debes poder explicar.}
\begin{itemize}[leftmargin=*]
  \item Por qué el Dummy tiene 66,1\% de accuracy y F1 = 0.
  \item Qué significa ``detecta 149 de 277''.
  \item Qué significa ``cuando alerta, acierta 74\%''.
  \item Diferencia entre ROC-AUC y PR-AUC.
  \item Que el umbral 0,5 es una decisión operativa, no una ley.
\end{itemize}

\begin{tcolorbox}[frase,title={Frase para el profesor}]
``No usamos accuracy como métrica principal porque el 66\% del test son negativos. El Dummy acierta el 66\% y no detecta nada. Por eso miramos F1, recall y PR-AUC.''
\end{tcolorbox}

\subsection{Bloque 5 --- Interpretabilidad y MILP (90 min)}

\paragraph{Estudia.}
\begin{itemize}[leftmargin=*]
  \item \code{fase 1 api/fase 2/reportes/importancia_variables_rf.png}.
  \item \code{fase 1 api/fase 2/src/optimizacion_milp.py}, líneas 78--128.
\end{itemize}

\paragraph{Debes poder explicar.}
\begin{itemize}[leftmargin=*]
  \item Top 5 de importancia y los bloques.
  \item Que Gini no es causalidad.
  \item La función objetivo del MILP.
  \item Las tres restricciones: máximo 4 brigadas, mínimo 2 con resguardos, distancia $\le$ 350 km.
  \item Por qué Pueblo Bello entra con 36,9\%.
  \item Que es \textbf{asignación}, no ruteo de vehículos.
\end{itemize}

\paragraph{Ejercicio.}
Escribe la función objetivo y las restricciones en una hoja, sin mirar el código. Luego resuelve un ejemplo de tres municipios por enumeración.

\subsection{Bloque 6 --- Demo y ensayo (60 min)}

\paragraph{Ejecuta.}
\begin{verbatim}
python "fase 2/src/probar_ia.py" "Valledupar"
python "fase 2/src/probar_ia.py" "Pueblo Bello"
python "fase 2/src/probar_ia.py" "Astrea"
python "fase 2/src/optimizacion_milp.py"
\end{verbatim}

\paragraph{Debes poder explicar.}
\begin{itemize}[leftmargin=*]
  \item Qué carga \code{probar_ia.py}: parquet + \code{modelo_gradient_boosting.joblib}.
  \item Por qué no reentrena.
  \item Por qué Valledupar da 0,9\% y Astrea 90,9\%.
  \item Por qué Pueblo Bello da 36,9\% y aun así entra al plan.
\end{itemize}

\paragraph{Ensayo.}
Graba un video de seis minutos y míralo. Repite sin notas.

\subsection{Checklist final}

\begin{itemize}[leftmargin=*,label=$\square$]
  \item Puedo dibujar el pipeline Fase~1 $\rightarrow$ dataset $\rightarrow$ Fase~2 $\rightarrow$ MILP $\rightarrow$ demo.
  \item Puedo escribir la fórmula del target sin mirar.
  \item Puedo explicar 74 columnas vs 27 features.
  \item Puedo explicar por qué el split es por municipio.
  \item Puedo recalcular precision, recall, F1 y accuracy de HGB.
  \item Puedo explicar el top 5 y los bloques de importancia.
  \item Puedo escribir la función objetivo y las restricciones del MILP.
  \item Puedo explicar Pueblo Bello sin titubear.
  \item Puedo correr la demo y explicar qué carga.
  \item Puedo decir tres limitaciones y tres mejoras.
  \item Puedo dar el guion de seis minutos sin notas.
\end{itemize}

% ============================================================
\section{Guion de seis minutos}
% ============================================================

\begin{table}[H]
\centering
\small
\begin{tabularx}{\textwidth}{@{}L{2.0cm} Y@{}}
\toprule
\textbf{Tiempo} & \textbf{Qué dices} \\
\midrule
0:00--0:30 &
``Las brigadas se asignan hoy con criterios empíricos. Nuestro objetivo es priorizar municipios con datos y optimizar su asignación.'' \\
0:30--2:00 &
``La Fase~1 construyó una tabla municipio--año: 1.122 municipios $\times$ 5 años, 74 columnas. De ahí salen 27 entradas y una salida: \code{prioridad_alta}, el tercio superior de la brecha de consultas frente a municipios pares.'' \\
2:00--3:00 &
``La Fase~2 compara cuatro modelos. HGB gana con F1 0,623, recall 0,538 y ROC-AUC 0,785. Detecta 149 de 277 municipios prioritarios y, cuando alerta, acierta 74\%.'' \\
3:00--4:00 &
``Las variables más importantes son NBI, perfil materno e infeccioso y población en resguardos. Es asociación, no causalidad.'' \\
4:00--5:00 &
``Con esa probabilidad, el MILP asigna cuatro brigadas desde Valledupar con restricciones de distancia y equidad. Resultado: Pueblo Bello, Becerril, Astrea y Chimichagua.'' \\
5:00--6:00 &
``Limitaciones: la etiqueta es relativa, los datos RIPS están agregados por sede del prestador y el modelo es transversal. Como trabajo futuro: validación temporal, calibración y ruteo real.'' \\
\bottomrule
\end{tabularx}
\end{table}

\subsection{Consejos de oralidad}

\begin{itemize}[leftmargin=*]
  \item No leas la slide: úsala como apoyo.
  \item Empieza por el problema, no por la tabla.
  \item Di los números redondeados: ``casi 78\% de accuracy'', ``detecta 149 de 277''.
  \item Si no sabes algo, di: ``no lo medimos; lo reconocería como limitación''.
  \item Cierra con limitaciones y trabajo futuro, no con un ``gracias'' vacío.
\end{itemize}

\subsection{Alternativa de contingencia: Pitch relámpago de tres minutos}
Si el docente Tonny Jiménez solicita ir directo al grano o recorta el tiempo de intervención, utiliza esta estructura condensada de tres minutos:

\begin{table}[H]
\centering
\small
\begin{tabularx}{\textwidth}{@{}L{2.0cm} Y@{}}
\toprule
\textbf{Tiempo} & \textbf{Mensaje clave (Qué decir sin titubear)} \\
\midrule
0:00--1:00 &
``Las brigadas médicas extramurales se despachan hoy por intuición política o reclamos puntuales. Diseñamos un sistema en dos etapas: una predictiva que detecta déficit asistencial relativo usando datos 100\% oficiales (RIPS, REPS, NBI DANE y resguardos), y una prescriptiva que despacha brigadas optimizando rutas y equidad constitucional.'' \\
1:00--2:00 &
``Construimos un panel de 5.610 observaciones (1.122 municipios $\times$ 5 años). Entrenamos cuatro clasificadores con split agrupado por municipio para evitar fuga de información temporal. El mejor fue \textbf{HistGradientBoosting} con un F1 de 0,6234, ROC-AUC de 0,7848 y una precisión del 74,13\%: 3 de cada 4 municipios marcados como prioritarios realmente carecen de atención suficiente.'' \\
2:00--3:00 &
``La probabilidad del modelo alimenta un MILP resuelto con HiGHS. Desde el Hospital Rosario Pumarejo en Valledupar, con 4 brigadas, el modelo asignó Pueblo Bello y Becerril por cuota ética de resguardos, y Astrea y Chimichagua por alto déficit asistencial. Contamos con una demo en consola (\code{probar_ia.py}) que predice en milisegundos sin reentrenar.'' \\
\bottomrule
\end{tabularx}
\end{table}

% ============================================================
\section{Banco de preguntas del profesor}
% ============================================================

\rowcolors{2}{gris!12}{white}
\begin{longtable}{@{}L{5.95cm} L{10.05cm}@{}}
\toprule
\textbf{Pregunta} & \textbf{Respuesta corta} \\
\midrule
\endhead
¿Qué es una fila? &
Un municipio en un año; 1.122 $\times$ 5 = 5.610. \\
¿Qué es \code{prioridad_alta}? &
Tercio superior de la brecha de consultas contra municipios del mismo tamaño y año. \\
¿Por qué 33\% de positivos? &
Por construcción: es el percentil 66 de la brecha. No es prevalencia real. \\
¿Hay fuga de datos? &
No usamos volumen, tasas, benchmark ni brecha. El split es por municipio. Salvedad: imputación de medianas sobre todo el dataset y CV no agrupada. \\
¿Por qué HGB? &
Mejor F1 (0,623) y PR-AUC (0,741); captura no linealidades y es robusto. \\
¿Qué significa ROC-AUC 0,785? &
Probabilidad de ordenar un positivo por encima de un negativo al azar. \\
¿Qué significa recall 0,538? &
Detecta 149 de 277 prioritarios; se le escapan 128. \\
¿Por qué no accuracy? &
El Dummy tiene 66,1\% de accuracy y F1 = 0; la clase negativa domina. \\
¿Por qué Pueblo Bello entra con 36,9\%? &
Por la cuota de equidad: mínimo dos municipios con resguardos. \\
¿Por qué MILP si el modelo ya ordena? &
El modelo puntúa; el MILP decide un portafolio factible con presupuesto, distancia y equidad. \\
¿Es ruteo? &
No: es asignación de brigadas a municipios. El ruteo de vehículos es trabajo futuro. \\
¿Qué mejorarías? &
GroupKFold por municipio, usar validación, imputar solo con train, calibrar probabilidades y validar con años externos. \\
¿Por qué 27 y no 74? &
Las 74 columnas incluyen identificadores, volumen, tasas y derivadas del target. Las 27 son contexto estructural y perfil de morbilidad. \\
¿Por qué solo cinco CIE-10? &
Materno, infeccioso, respiratorio, digestivo y contacto con servicios son los más relevantes para brigadas; los otros 21 añaden ruido. \\
¿Usaron el set de validación? &
La selección fue con CV=3 sobre train; validación queda para calibrar el umbral y la siguiente iteración. \\
¿Cómo funciona la demo? &
Carga el parquet y el modelo HGB ya entrenados, filtra el municipio en 2021 y muestra probabilidad y recomendación. No reentrena. \\
¿Por qué $\ln(V_i + 1)$ en el MILP y no $V_i$ lineal? &
Modela rendimientos decrecientes; evita que un municipio gigante monopolice las brigadas y premia la equidad distributiva. \\
¿Cómo definieron $\lambda = 0{,}005$? &
Pondera el costo de viaje frente al impacto. A 200 km penaliza 1,0 punto al score; permite viajes lejanos solo si la vulnerabilidad es alta. \\
¿Por qué Haversine y no rutas viales? &
Haversine es determinista, autónoma y reproducible sin depender de servicios web ni costos de API; en Fase 3 se integrará matriz OSRM. \\
¿Qué pasaría si bajan el umbral a 0,35? &
El recall aumentaría por encima del 75\%, capturando más municipios desatendidos, a cambio de aceptar más falsas alarmas (menor precisión). \\
¿Cómo se calculó la población vulnerable $V_i$? &
Suma directa de población rural dispersa (Censo DANE 2018 / PPED) más población censada en resguardos indígenas (MinInterior). \\
¿Por qué imputaron medianas en CIE-10? &
Menos del 0,5\% de registros presentaban ausencia en diagnósticos secundarios; la mediana previene ceros espurios y distorsiones. \\
¿Qué solver resuelve el MILP? &
HiGHS mediante \code{scipy.optimize.linprog}, solver de programación lineal entera de alto desempeño basado en simplex dual paralelizado. \\
¿El sistema se puede aplicar a otros departamentos? &
Sí; el panel contiene los 1.122 municipios de Colombia y basta con indicar las coordenadas del hospital base departamental. \\
¿Por qué no incluyeron la distancia dentro del modelo predictivo? &
Porque la necesidad de salud es intrínseca a la población del municipio; la distancia a la base es un costo logístico prescriptivo, no un síntoma clínico. \\
\bottomrule
\end{longtable}

% ============================================================
\section{Errores que no debes decir}
% ============================================================

\begin{table}[H]
\centering
\small
\begin{tabularx}{\textwidth}{@{}L{7.0cm} Y@{}}
\toprule
\textbf{No digas} & \textbf{Di} \\
\midrule
``El modelo acierta el 78\% de los prioritarios'' &
``Tiene 77,97\% de accuracy global y detecta el 53,8\% de los prioritarios.'' \\
``18M de atenciones'' &
``1.800 millones de atenciones.'' \\
``1.122 municipios entrenados'' &
``785 municipios en train; 1.122 es el total del dataset.'' \\
``Cero data leakage'' &
``Sin fuga directa; hay salvedades metodológicas que reconocemos.'' \\
``Dummy azar'' &
``Dummy \code{most_frequent}: predice siempre la clase mayoritaria.'' \\
``El modelo decide las brigadas'' &
``El modelo predice la probabilidad de déficit asistencial; el MILP decide la asignación óptima de brigadas.'' \\
``Estas variables causan la prioridad'' &
``Estas variables son las más asociadas en el ensamble; no implican relación causal.'' \\
``El 33\% de Colombia tiene déficit crítico'' &
``El 33\% es el tercio superior de una brecha relativa construida frente a pares poblacionales.'' \\
``Hacemos ruteo de vehículos'' &
``Hacemos asignación y despacho de portafolio; el ruteo de vehículos es trabajo de la Fase 3.'' \\
``Usamos el set de validación para elegir hiperparámetros'' &
``La selección fue con CV=3 sobre train; validación queda reservada para calibrar umbrales y la siguiente iteración.'' \\
``Excluimos RIPS por fuga'' &
``Excluimos volumen y tasas directas de consultas; conservamos los porcentajes diagnósticos CIE-10 como perfil epidemiológico.'' \\
``El MILP calcula las mejores carreteras'' &
``El MILP asigna municipios minimizando distancia ortodrómica Haversine; no optimiza rutas sobre la red vial terrestre.'' \\
``La precisión o el recall son malos'' &
``La precisión es sólida (74,13\%: 3 de cada 4 alertas son reales); el recall moderado (53,79\%) refleja un umbral conservador de 0,5.'' \\
\bottomrule
\end{tabularx}
\end{table}

% ============================================================
\section{Auditoría y preparación diapositiva a diapositiva (Pitch Deck)}
% ============================================================

A continuación se audita el contenido de las 13 diapositivas de la presentación oficial de la Fase~2 (\code{IA - SEGUNDA FASE.pdf}), indicando qué mensajes transmitir, qué trampas evitar y qué preguntas esperar del jurado en cada lámina.

\subsection{Lámina 1: Portada Institucional}
\begin{itemize}[leftmargin=*]
  \item \textbf{Contenido visual:} Título oficial del proyecto, participantes (Jorman Noriega, Julio Ríos, Gabriel Venegas), docente Tonny Jiménez, Universidad Popular del Cesar, Valledupar, 2026.
  \item \textbf{Qué decir:} ``Presentamos el sistema inteligente de priorización y optimización espacial de brigadas de salud extramurales para comunidades indígenas y rurales dispersas, correspondiente a la Fase~2 del proyecto de aula de Inteligencia Artificial.''
  \item \textbf{Cuidado:} Mantener seguridad, postura erguida y modular la voz desde el primer segundo.
\end{itemize}

\subsection{Lámina 2: Planteamiento del Problema}
\begin{itemize}[leftmargin=*]
  \item \textbf{Contenido visual:} Solicitudes informales, presupuestos fijos, 1. Inequidad en cobertura, 2. Uso ineficiente de recursos, 3. Falta de predictibilidad epidemiológica.
  \item \textbf{Qué decir:} ``El despacho actual de brigadas en Colombia adolece de tres fallas estructurales: se basa en solicitudes informales o presiones coyunturales, genera gasto ineficiente en traslados y no anticipa el déficit crónico en comunidades dispersas.''
  \item \textbf{Cuidado:} No culpar a una entidad particular; presentar el problema como una oportunidad de mejora cuantitativa en la gestión pública.
\end{itemize}

\subsection{Lámina 3: Enfoque Metodológico}
\begin{itemize}[leftmargin=*]
  \item \textbf{Contenido visual:} Aplicada, Cuantitativa, Correlacional-Explicativa, CRISP-DM adaptada a salud.
  \item \textbf{Qué decir:} ``Es una investigación aplicada y cuantitativa. Analizamos la correlación entre capacidad instalada, pobreza estructural y demanda médica para alimentar un pipeline basado en CRISP-DM.''
  \item \textbf{Trampa a evitar:} Si la diapositiva dice ``Relaciones Causa-Efecto'', aclara inmediatamente: \textbf{``Analizamos correlaciones y patrones predictivos; no establecemos relaciones de causalidad médica directa, pues trabajamos con datos observacionales agregados.''}
\end{itemize}

\subsection{Lámina 4: Materia Prima (Datos Abiertos 100\% Oficiales)}
\begin{itemize}[leftmargin=*]
  \item \textbf{Contenido visual:} DIVIPOLA, RIPS, REPS, Resguardos MinInterior, Pobreza NBI Censo 2018, Proyecciones PPED DANE.
  \item \textbf{Qué decir:} ``Integramos seis fuentes oficiales del Estado colombiano para conformar un panel analítico municipal de 5.610 observaciones y 74 variables con centroides GPS 100\% validados.''
  \item \textbf{Correcciones obligatorias en la exposición:}
\end{itemize}

\begin{table}[H]
\centering
\small
\begin{tabularx}{\textwidth}{@{}L{4.2cm} Y@{}}
\toprule
\textbf{Dice la lámina} & \textbf{Debes decir oralmente} \\
\midrule
``782 municipios'' &
``1.122 municipios totales a nivel nacional (785 conformaron el conjunto de entrenamiento).'' \\
``18 mil millones de atenciones'' &
``1.800 millones de atenciones registradas en RIPS entre 2017 y 2021.'' \\
``14 dimensiones de vulnerabilidad'' &
``14 indicadores de NBI construidos sobre las 5 necesidades básicas del DANE.'' \\
``causas CIE-10'' &
``perfil de diagnósticos CIE-10 (RIPS registra diagnósticos atendidos, no causas básicas).'' \\
``delimitación de resguardos'' &
``conteo de resguardos y población censada en resguardos (datos tabulares, no polígonos).'' \\
``DANE / MinSalud'' para DIVIPOLA &
``DANE: DIVIPOLA es la codificación oficial de la autoridad estadística nacional.'' \\
Ligaduras rotas (\textit{geográ cos}) &
Pronunciar correctamente: ``geográficos oficiales''. \\
\bottomrule
\end{tabularx}
\end{table}

\subsection{Lámina 5: El Secreto del Target: Equidad Territorial}
\begin{itemize}[leftmargin=*]
  \item \textbf{Contenido visual:} Fórmula de tasa de consultas por 1.000 hab, benchmark por quintil poblacional (\code{banda_pob}), brecha asistencial relativa, umbral percentil 66 ($0{,}337$).
  \item \textbf{Qué decir:} ``No comparamos un caserío rural contra Bogotá o Medellín. Comparamos cada municipio contra sus pares de tamaño similar en el mismo año. Aislamos los prestadores de referencia regional (\textit{hubs}) y etiquetamos como prioridad alta al tercio con mayor brecha (percentil 66).''
  \item \textbf{Pregunta típica:} ``¿Por qué hay 33,3\% de positivos?'' $\rightarrow$ ``Por definición matemática del percentil 66, no por prevalencia epidemiológica.''
\end{itemize}

\subsection{Lámina 6: Partición 70 / 15 / 15 Agrupada por Municipio}
\begin{itemize}[leftmargin=*]
  \item \textbf{Contenido visual:} Train 70\% (785 municipios, 3.770 válidos), Val 15\% (168 municipios, 807 válidos), Test 15\% (169 municipios, 817 válidos).
  \item \textbf{Qué decir:} ``La partición se agrupó por código municipal (\code{cod_mpio}). Ningún municipio del conjunto de prueba ciego fue visto jamás por el modelo durante el entrenamiento, garantizando una evaluación rigurosa de generalización territorial.''
  \item \textbf{Honestidad académica:} ``Reconocemos dos salvedades técnicas: la imputación de medianas fue global y la validación cruzada de GridSearchCV no agrupó por municipio; ambas se corrigen en la Fase~3.''
\end{itemize}

\subsection{Lámina 7: Torneo de Modelos: HistGradientBoosting es el Ganador}
\begin{itemize}[leftmargin=*]
  \item \textbf{Contenido visual:} Tabla comparativa de los cuatro clasificadores (Accuracy, Precision, Recall, F1, ROC-AUC, PR-AUC).
  \item \textbf{Qué decir:} ``HistGradientBoosting superó ampliamente a la Regresión Logística y al Random Forest, alcanzando un F1 de 0,6234, ROC-AUC de 0,7848 y una precisión del 74,13\%. En el Dummy baseline demostramos que una accuracy del 66,1\% es ficticia porque no detecta ningún caso crítico.''
  \item \textbf{Dato a enfatizar:} ``El modelo acierta 3 de cada 4 alertas de prioridad (precisión 74,13\%), optimizando el uso de recursos fiscales limitados.''
\end{itemize}

\subsection{Lámina 8: Importancia de Variables (Gini)}
\begin{itemize}[leftmargin=*]
  \item \textbf{Contenido visual:} Gráfica de importancia del Random Forest; dominancia de diagnósticos maternos (\code{pct_cap_O}), población en resguardos y NBI rural.
  \item \textbf{Qué decir:} ``El ensamble prioriza municipios combinando demanda obstétrica e infecciosa con pobreza estructural: el bloque de NBI concentra el 42\% de la importancia total y el perfil CIE-10 el 26\%.''
  \item \textbf{Aclaración clave:} ``Gini mide reducción de impureza en las divisiones de los árboles; indica asociación estadística predictiva, no causalidad médica.''
\end{itemize}

\subsection{Lámina 9: Inteligencia Prescriptiva (MILP)}
\begin{itemize}[leftmargin=*]
  \item \textbf{Contenido visual:} Función objetivo con $\ln(V_i + 1) - \lambda d_i$ y restricciones de flota ($K=4$), cuota indígena ($M=2$) y radio operativo ($D_{\max}=350$ km).
  \item \textbf{Qué decir:} ``La inteligencia prescriptiva cierra la brecha entre la predicción y la acción logística. Formulamos un MILP resuelto con HiGHS que maximiza la cobertura de vulnerabilidad penalizando la distancia de traslado desde la base operativa.''
  \item \textbf{Cuidado:} Dejar claro que es \textbf{asignación óptima de portafolio}, no ruteo de caminos.
\end{itemize}

\subsection{Lámina 10: Demostración Operativa en el Cesar}
\begin{itemize}[leftmargin=*]
  \item \textbf{Contenido visual:} Tabla de despacho desde el Hospital Rosario Pumarejo de Valledupar: Pueblo Bello (36,9\% prob, 2 resguardos, 86,6\% NBI rural), Becerril (45,8\%, 2 resguardos), Astrea (90,9\% prob), Chimichagua (75,0\% prob).
  \item \textbf{Qué decir:} ``El solver cumplió con exactitud la cuota ética asignando Pueblo Bello y Becerril en la Sierra Nevada y el Perijá, y asignó las dos brigadas restantes a Astrea y Chimichagua por su crítico abandono ambulatorio.''
\end{itemize}

\subsection{Lámina 11: Demostración Funcional en Consola}
\begin{itemize}[leftmargin=*]
  \item \textbf{Contenido visual:} Captura de ejecución de \code{probar_ia.py} y \code{optimizacion_milp.py}.
  \item \textbf{Qué decir:} ``Diseñamos una herramienta de consola reproducible que recupera el vector municipal de 27 variables y emite la predicción y el plan en milisegundos sin necesidad de reentrenar.''
\end{itemize}

\subsection{Lámina 12: Limitaciones Reconocidas del Enfoque}
\begin{itemize}[leftmargin=*]
  \item \textbf{Contenido visual:} Cuatro limitaciones técnicas bien definidas.
  \item \textbf{Qué decir:} ``Reconocemos con rigor científico nuestras limitaciones: RIPS agrega consultas por sede del prestador y no por residencia del paciente; la distancia es ortodrómica Haversine; y la validación no se usó para early stopping ni calibración de umbrales.''
\end{itemize}

\subsection{Lámina 13: Hoja de Ruta para Fase 3 y Conclusiones}
\begin{itemize}[leftmargin=*]
  \item \textbf{Contenido visual:} Aplicación web en Streamlit con mapas Folium, calibración de umbrales con validación, ruteo OSRM por carreteras y video pitch de 3 minutos.
  \item \textbf{Qué decir:} ``La Fase~2 dejó el núcleo matemático y algorítmico completamente probado. En la Fase~3 integraremos la interfaz interactiva con mapas satelitales para que las Secretarías de Salud tomen decisiones basadas en datos en tiempo real.''
\end{itemize}

% ============================================================
\section{Mapa de evidencia y comandos}
% ============================================================

\subsection{Dónde verificar cada afirmación}

\begin{table}[H]
\centering
\small
\begin{tabularx}{\textwidth}{@{}L{6.2cm} Y@{}}
\toprule
\textbf{Afirmación} & \textbf{Dónde verificarla} \\
\midrule
5.610 filas, 74 columnas, 1.122 municipios &
\code{fase 1 api/data/modelo/municipio_anio.parquet} y \code{fase 1 api/fase 1/docs/diccionario_datos.md}. \\
Target y split &
\code{fase 1 api/fase 1/src/build_tabla_modelo.py}, líneas 38--56 y 120--154. \\
27 features &
\code{fase 1 api/fase 2/src/modelado.py}, líneas 36--69. \\
Métricas y matriz de confusión &
\code{fase 1 api/fase 2/reportes/metricas_comparacion.csv} y \code{matriz_confusion_modelos.png}. \\
Importancia de variables &
\code{fase 1 api/fase 2/reportes/importancia_variables_rf.png}. \\
MILP &
\code{fase 1 api/fase 2/src/optimizacion_milp.py}, líneas 34--128. \\
Demo &
\code{fase 1 api/fase 2/src/probar_ia.py}. \\
\bottomrule
\end{tabularx}
\end{table}

\subsection{Comandos para reproducir la demo}

\begin{verbatim}
cd "fase 1 api"
python "fase 2/src/probar_ia.py" "Valledupar"
python "fase 2/src/probar_ia.py" "Pueblo Bello"
python "fase 2/src/optimizacion_milp.py"
\end{verbatim}

\subsection{Cómo compilar este documento}

\begin{verbatim}
cd "guia_estudio"
pdflatex guia_estudio_fase2.tex
pdflatex guia_estudio_fase2.tex
\end{verbatim}

% ============================================================
\appendix
\section{Glosario}
% ============================================================

\begin{description}[leftmargin=*,style=nextline,itemsep=4pt]
  \item[Municipio--año] Una fila del dataset: un municipio en un año concreto.
  \item[Hub] Prestador de referencia regional cuya tasa de consultas es al menos 3 veces la mediana de su banda. Se excluye del entrenamiento.
  \item[NBI] Necesidades Básicas Insatisfechas. Mide pobreza multidimensional con cinco componentes: vivienda, servicios, hacinamiento, inasistencia escolar y dependencia económica.
  \item[CIE-10] Clasificación Internacional de Enfermedades, décima revisión. En el dataset se usan porcentajes por capítulo diagnóstico.
  \item[Gini importance] Reducción media de impureza que aporta una variable en un bosque aleatorio. No es causalidad.
  \item[MILP] Programación Lineal Entera Mixta. Optimiza una función objetivo con variables binarias y restricciones lineales.
  \item[Haversine] Distancia geodésica entre dos puntos de la Tierra a partir de latitud y longitud.
  \item[Precision] De los municipios que el modelo marca como prioritarios, qué proporción realmente lo son.
  \item[Recall] De los municipios que realmente son prioritarios, qué proporción detecta el modelo.
  \item[F1] Media armónica entre precision y recall.
  \item[ROC-AUC] Probabilidad de que el modelo ordene un positivo por encima de un negativo al azar.
  \item[PR-AUC] Área bajo la curva precision--recall; útil cuando la clase positiva es minoritaria.
  \item[Data leakage] Uso de información que no estará disponible al predecir, o que contiene la respuesta.
  \item[Split] Partición train/validación/test. En este proyecto se hace por municipio para evitar fuga.
\end{description}

\section{Checklist de archivos}

\begin{itemize}[leftmargin=*,label=$\square$]
  \item \code{municipio_anio.parquet} existe y tiene 5.610 filas.
  \item \code{diccionario_datos.md} está actualizado.
  \item \code{modelado.py} contiene las 27 features.
  \item \code{metricas_comparacion.csv} tiene los cuatro modelos.
  \item \code{matriz_confusion_modelos.png} y \code{curva_roc_comparacion.png} existen.
  \item \code{optimizacion_milp.py} corre y genera \code{asignacion_brigadas_optima.csv}.
  \item \code{probar_ia.py} responde para Valledupar, Pueblo Bello y Astrea.
  \item \code{documento.pdf} de la Fase~2 tiene 12 páginas.
  \item Esta guía compila sin errores.
\end{itemize}

\vfill
\begin{center}
\small\color{grisoscuro}
Guía de estudio y defensa --- Fase 2\\
Universidad Popular del Cesar --- Ingeniería de Sistemas e Informática\\
2026
\end{center}

\end{document}
